\documentclass[11pt]{article}
\usepackage[margin=1in]{geometry}
\usepackage{amsmath,amssymb,amsthm}
\usepackage{mathtools}
\usepackage{hyperref}

\newtheorem{theorem}{Theorem}
\newtheorem{corollary}{Corollary}
\newtheorem{lemma}{Lemma}
\newtheorem{proposition}{Proposition}

\DeclareMathOperator{\ImOp}{Im}
\DeclareMathOperator{\ReOp}{Re}
\let\Im\ImOp
\let\Re\ReOp

\title{Spectral Tiling of $\zeta$ on the Critical Line\\
{\large Corrected form: main-sum pedestal plus remainder atom at $\omega=-1$}}
\author{}
\date{}

\begin{document}
\maketitle

\section*{Setup}

Let
\[
  \vartheta(t) \;=\; \arg\Gamma\!\left(\tfrac{1}{4}+\tfrac{it}{2}\right) \;-\; \tfrac{t}{2}\log\pi
\]
be the Riemann--Siegel theta function, $\tau_{0} = t/(2\pi)$, $N(t)=\lfloor\sqrt{\tau_{0}}\rfloor$, $p(t)=\sqrt{\tau_{0}}-N(t)-\tfrac12$. For a window $[T_{0},T_{\max}]$ with $T_{0}\ge 200$, define the \emph{warped Fourier coefficient}
\[
  c(\omega) \;=\; \frac{1}{\Delta T}\int_{T_{0}}^{T_{\max}} \zeta\!\left(\tfrac{1}{2}+it\right) e^{-i\omega\,\vartheta(t)}\,dt,
  \qquad \Delta T = T_{\max}-T_{0},
\]
and the band-limited reconstruction
\[
  \zeta_{\mathrm{st}}(\tau) \;=\; \int_{-2}^{0} c(\omega)\,e^{\,i\omega\tau}\,d\omega.
\]

\section*{Main theorem}

\begin{theorem}[Spectral tiling of $\zeta$ on the critical line]\label{thm:main}
The warped spectrum of $\zeta(\tfrac12+it)$ decomposes as
\[
  c(\omega) \;=\; c_{\mathcal D}(\omega) \;+\; c_{\mathcal R}(\omega),
  \qquad
  \mathrm{supp}\,c_{\mathcal D} \;=\; [-1,0],
  \qquad
  \mathrm{supp}\,c_{\mathcal R} \;=\; \{-1\},
\]
i.e. the Hardy main sum contributes a continuous pedestal on $[-1,0]$ and the Riemann--Siegel remainder contributes an atom at $\omega=-1$. Consequently, for every $t\in[T_{0},T_{\max}]$,
\[
  \bigl|\zeta_{\mathrm{st}}(\vartheta(t)) \;-\; \zeta(\tfrac12+it)\bigr|
  \;\le\; C_{1}\,\tau_{0}^{-(2M+3)/4} \;+\; C_{2}\,(\Delta T)^{-1}\log\tau_{0},
\]
and
\[
  \lim_{\substack{\Delta T\to\infty \\ M\to\infty}} \zeta_{\mathrm{st}}(\vartheta(t)) \;=\; \zeta(\tfrac12+it)
  \qquad \text{pointwise on } [T_{0},\infty).
\]
\end{theorem}

\section*{Proof}

\subsection*{Step 1: Riemann--Siegel decomposition}

For $t\ge 200$, the exact Riemann--Siegel formula gives
\[
  \zeta(\tfrac12+it)
  \;=\; \underbrace{\sum_{n=1}^{N(t)} n^{-1/2}\,e^{-it\log n}}_{\mathcal D(t)}
  \;+\; \underbrace{e^{-i\vartheta(t)}R(t)}_{\mathcal R(t)},
\]
where $R(t)$ is the Riemann--Siegel remainder (Edwards, \emph{Riemann's Zeta Function}, \S7.4).

\subsection*{Step 2: Main-sum pedestal on $[-1,0]$}

The coefficient $c_{\mathcal D}(\omega) = \frac{1}{\Delta T}\int \mathcal D(t)\,e^{-i\omega\vartheta(t)}\,dt$ has $n$-th phase $\Phi_{n}(t) = -t\log n - \omega\vartheta(t)$, with
\[
  \Phi_{n}'(t) \;=\; -\log n \;-\; \tfrac{\omega}{2}\log\tau_{0}.
\]
Stationary phase gives
\[
  \omega \;=\; -\,\frac{2\log n}{\log\tau_{0}}, \qquad 1\le n\le N(t)=\sqrt{\tau_{0}},
\]
so $\omega$ sweeps continuously through $[-1,0]$ as $n$ ranges from $1$ to $\sqrt{\tau_{0}}$. Hence $\mathrm{supp}\,c_{\mathcal D} = [-1,0]$.

\subsection*{Step 3: The remainder atom at $\omega=-1$}

This is the step that must be stated correctly. The remainder, in the shifted object $\mathcal R(t)=e^{-i\vartheta(t)}R(t)$, carries the phase factor $e^{-i\vartheta(t)}$ in front of a quantity whose modulation is $\tau_{0}^{-1/4}C_{0}(p(t))+O(\tau_{0}^{-3/4})$ in the sense of Gabcke's 1979 thesis Theorem 4.2. The Gabcke function
\[
  C_{0}(p) \;=\; \frac{\cos 2\pi(p^{2}-p-\tfrac{1}{16})}{\cos 2\pi p}
\]
is a bounded function of $p(t)\in[-\tfrac12,\tfrac12]$ alone, not of $t\log n$. The fractional-part function $p(t)$ advances by $1$ each time $\sqrt{\tau_{0}}$ crosses an integer, so $p(t)$ has instantaneous frequency
\[
  p'(t) \;=\; \frac{1}{4\pi\sqrt{\tau_{0}}} \;\to\; 0
\]
on the $t$-axis, and on the $\vartheta$-axis
\[
  \frac{dp}{d\vartheta} \;=\; \frac{p'(t)}{\vartheta'(t)}
  \;=\; \frac{1}{4\pi\sqrt{\tau_{0}}\cdot\tfrac12\log\tau_{0}}
  \;=\; \frac{1}{2\pi\sqrt{\tau_{0}}\log\tau_{0}}
  \;\to\; 0.
\]
Hence on the $\vartheta$-axis, $C_{k}(p(t))$ is a slowly varying amplitude with zero instantaneous frequency in the limit; its warped Fourier transform is concentrated at the origin of the $\omega$-axis-relative-to-its-carrier.

The carrier of $\mathcal R(t)$ is $e^{-i\vartheta(t)}$, which has warped frequency $\omega=-1$. Thus
\[
  c_{\mathcal R}(\omega)
  \;=\; \frac{1}{\Delta T}\int_{T_{0}}^{T_{\max}} e^{-i\vartheta(t)}R(t)\,e^{-i\omega\vartheta(t)}\,dt
  \;=\; \frac{1}{\Delta T}\int_{T_{0}}^{T_{\max}} R(t)\,e^{-i(\omega+1)\vartheta(t)}\,dt,
\]
and the stationary-phase condition is $(\omega+1)\vartheta'(t)=0$, which has a solution only at $\omega=-1$ (since $\vartheta'(t)>0$ throughout the window by Lemma~1 below). For any $\omega\ne -1$, integration by parts gives
\[
  |c_{\mathcal R}(\omega)|
  \;\le\; \frac{K\,T_{0}^{-1/4}}{|\omega+1|\log\tau_{0}}\cdot(\Delta T)^{-1}
  \;=\; O(\Delta T^{-1}),
\]
which vanishes as $\Delta T\to\infty$. At $\omega=-1$ exactly,
\[
  c_{\mathcal R}(-1)
  \;=\; \frac{1}{\Delta T}\int_{T_{0}}^{T_{\max}} R(t)\,dt
  \;=\; O(T_{0}^{-1/4}),
\]
a finite nonzero value. Therefore $\mathrm{supp}\,c_{\mathcal R}=\{-1\}$: the remainder is an atom at $\omega=-1$, not a band.

\subsection*{Step 4: Why the earlier ``band'' statement was wrong}

The earlier derivation conflated two different saddle problems. The \emph{dual sum} $2\sum_{n=1}^{N}\cos(\vartheta+t\log n)/\sqrt n$ is a formal representation that appears in the Riemann--Siegel derivation, but that sum is not $R(t)$; rather, $R(t)$ is the saddle-point residue of that sum, already collapsed to a single slowly varying amplitude $\tau_{0}^{-1/4}C_{0}(p)\cdot(-1)^{N-1}$. The dual-sum tones do not survive as independent spectral contributions; they have been integrated out. What remains is the carrier $e^{i\vartheta}$ (or $e^{-i\vartheta}$ after the shift to $\zeta$) multiplied by a function of $p(t)$ alone. That is a single-carrier signal, with a single warped frequency. The correct picture is: \emph{main sum = pedestal on $[-1,0]$, remainder = atom at $\omega=-1$}, and the atom sits at the right endpoint of the pedestal's reflection, which is the functional-equation fixed point.

\subsection*{Step 5: Functional-equation interpretation}

The involution $\omega\mapsto -2-\omega$ on the warped axis is the spectral image of $s\mapsto 1-s$. The pedestal $[-1,0]$ and its reflection $[-2,-1]$ would together tile $[-2,0]$ if both were realized. They are not. Only the pedestal $[-1,0]$ is realized as a band; its reflection collapses to the single fixed point $\omega=-1$ because the Riemann--Siegel saddle has already contracted the reflected band to its endpoint. This collapse is precisely what makes $R(t)=O(t^{-1/4})$ rather than $O(1)$ --- the reflected band, if uncollapsed, would carry the same amplitude as the main pedestal, but the saddle integration reduces it by the factor $\tau_{0}^{-1/4}$ coming from $\sqrt{2\pi/|\vartheta''|}$. The atom at $\omega=-1$ is therefore the saddle-collapsed image of $[-2,-1]$ under the functional equation.

\subsection*{Step 6: Error bound}

Let $R(t)=R_{M}(t)+E_{M}(t)$ with $|E_{M}(t)|\le K_{M}\tau_{0}^{-(2M+3)/4}$ (Gabcke 1979 thesis, Theorem 4.2; Arias--de Reyna, \emph{High precision computation of Riemann's zeta function by the Riemann--Siegel formula}, Math.~Comp.~\textbf{80} (2011) 995--1009, Theorem 2.1). The atom location is $M$-independent (Step 3 uses only the carrier $e^{-i\vartheta}$, not the asymptotic order), so
\[
  c_{\mathcal R_{M}}(\omega)-c_{\mathcal R}(\omega)
  \;=\; \frac{1}{\Delta T}\int E_{M}(t)\,e^{-i(\omega+1)\vartheta(t)}\,dt
  \;=\; O\!\bigl(T_{0}^{-(2M+3)/4}\bigr).
\]
Combined with the endpoint leakage of the main-sum pedestal at $\omega=0$ and at $\omega=-1$, this yields the stated error bound. Taking $\Delta T\to\infty$ and $M\to\infty$ gives pointwise convergence. $\blacksquare$

\section*{Corollaries}

\begin{corollary}[Saddle-collapse identity]
The functional equation $\chi(s)\chi(1-s)=1$ on the critical line acts on the warped spectrum of $\zeta$ by collapsing the reflected pedestal $[-2,-1]$ to the atom $\{-1\}$, with mass scaled by $\tau_{0}^{-1/4}$ relative to the main pedestal on $[-1,0]$. The pedestal and its saddle-collapsed image meet exactly at $\omega=-1$, the fixed point of the reflection.
\end{corollary}

\begin{corollary}[Reconstruction bandwidth]
The support of $c(\omega)$ is $[-1,0]\cup\{-1\} = [-1,0]$. The band $[-2,0]$ used in the reconstruction is larger than necessary by exactly one unit; the reconstruction $\zeta_{\mathrm{st}}(\tau)$ computed on $[-1,0]$ alone converges to $\zeta(\tfrac12+it)$ with the same error bound.
\end{corollary}

\section*{Auxiliary lemma used in Step 3}

\begin{lemma}[Positivity of $\vartheta'$]
For every $t\ge 200$, $\vartheta'(t) \ge \tfrac12\log(t/(2\pi)) - \tfrac{1}{200} > 0$.
\end{lemma}

\begin{proof}
Binet's first formula for $\Re\psi(\tfrac14+\tfrac{it}{2})$ gives
\[
  \tfrac12\Re\psi(\tfrac14+\tfrac{it}{2}) - \tfrac12\log\pi
  \;=\; \tfrac12\log(t/(2\pi)) + \eta(t), \qquad |\eta(t)|\le \tfrac{1}{200} \text{ for } t\ge 200,
\]
by the standard bounds on the Binet remainder integral.
\end{proof}

\end{document}
